\documentclass{osucourse}
\course{5576H}

%% Not in the university's course data, so these come
%% from the published sheet this file was imported from.
\SetCourseField{title}{Honors Number Theory}
\SetCourseField{credits}{5 credits}
\SetCourseField{description}{Elementary analytic and algebraic number theory, tracing its unifying role in the development of mathematics through history.}
\SetCourseField{prereq}{C or better in 4182H, or in both 2182H and 3345; or credit for 264H, or for both 263H and 345; or permission of department.}

\begin{document}

\CatalogDescription
\PrerequisitesAndExclusions

\begin{Purpose}
The intention of this course is to present number theory, the "Queen of Mathematics" through its historical development. Being one of the oldest mathematical disciplines, number theory, in the course of its history, both benefited from and contributed to such major mathematical areas as geometry, algebra and analysis. These courses will be especially beneficial for honor students planning to pursue careers in mathematics, physics, computer science and education, but may be of interest to engineering students as well.
\end{Purpose}

\begin{Text}
Vary, for example:  \emph{An Introduction to the Theory of Numbers}, 6\textsuperscript{th} edition, by Hardy, Wright, Heath \& Brown, published by Oxford, ISBN: 9780199219865.  An Introduction to the Theory of Numbers, I. Niven, H.S. Zukkerman, H.L. Montgomery  \emph{Number Theory: An Introduction to Mathematics, Parts A and B}, by William A. Coppel, Springer-Velag.
\end{Text}

\begin{TopicsList}
\item Review of Egyptian and Mesopotamian Mathematics. Greek tradition. Three classical Greek problems (cube doubling, angle trisection, circle quadrature).
\item Famous irrationalities.
\item Continued fractions and applications thereof (quadratic surds, Pell’s equation, Diophantine approximations, etc.)
\item More on diophantine approximation. Algebraic numbers. Liouville numbers. A glimpse into the Thue-Siegel-Roth Theorem.
\item Uniform distribution modulo one. Weyl criterion. Some important sequences. Pisot-Vijayaraghavan numbers. Formulation and discussion of Margulis’ solution of Oppenheimer’s conjecture.
\item Normal numbers. Champernoun’s example. Almost every number is normal. Levy-Khinchine Theorem on normality of continued fractions.
\item Infinitude of primes. Euler’s identity. Chebyshev’s Theorem. Bertrand’s Postulate. Dirichlet’s Theorem on primes in progressions. Average rate of growth of classical number-theoretical functions.
\item Finite fields. Wedderburn’s Theorem. Applications: Latin Squares and Cryptography.
\item Quadratic reciprocity.
\item Pythagorean triangles. Representation of integers as sums of squares. Quaternions, Cayley’s octavas. Hurwitz’ Theorem. Minkowski’s geometry of numbers.
\item p-adic numbers, their construction and axiomatic characterization (Ostrowski’s Theorem). Minkowski-Hasse principle.
\item Fermat’s last theorem. Some easy cases. A glimpse into modern developments (elliptic curves, Mordell-Weil Theorem, etc.).
\end{TopicsList}

\end{document}
